\documentclass[10pt,a4paper]{amsart}
\usepackage{amsmath,amssymb,amsthm,graphicx}
% Retained preamble for the Stacks Project chapter {CH} (shared project preamble).
% The chapter file pinned-original-source.tex only carries \input{{preamble}} on its first line,
% so no preamble text lives inside the chapter itself.  The Stacks Project uses one single shared
% preamble for every chapter: the snapshot stages it three times (sources/stacks-chow/preamble.tex,
% sources/stacks-etale-cohomology/preamble.tex and focused-seeds/AGP/sources/stacks-snapshot-20261011/
% preamble.tex) and those copies are identical in size.  The block below is copied mechanically,
% line for line, from sources/stacks-chow/preamble.tex lines 14-18 and 154-217: the Xy-pic setup and
% the native theorem declarations together with the mathematical macros.  Nothing was added, removed
% or redefined; the wrapper class and the amsmath/amssymb/amsthm/graphicx base packages are supplied
% by the assembler.
% For commutative diagrams we use Xy-pic
% Chapter: stacks-cohomology.tex -- Cohomology of Sheaves.
\usepackage[all]{xy}

% We use 2cell for 2-commutative diagrams.
\xyoption{2cell}

% Theorem environments.
%
\theoremstyle{plain}
\newtheorem{theorem}[subsection]{Theorem}
\newtheorem{proposition}[subsection]{Proposition}
\newtheorem{lemma}[subsection]{Lemma}

\theoremstyle{definition}
\newtheorem{definition}[subsection]{Definition}
\newtheorem{example}[subsection]{Example}
\newtheorem{exercise}[subsection]{Exercise}
\newtheorem{situation}[subsection]{Situation}

\theoremstyle{remark}
\newtheorem{remark}[subsection]{Remark}
\newtheorem{remarks}[subsection]{Remarks}

\numberwithin{equation}{subsection}

% Macros
%
\def\lim{\mathop{\mathrm{lim}}\nolimits}
\def\colim{\mathop{\mathrm{colim}}\nolimits}
\def\Spec{\mathop{\mathrm{Spec}}}
\def\Hom{\mathop{\mathrm{Hom}}\nolimits}
\def\Ext{\mathop{\mathrm{Ext}}\nolimits}
\def\SheafHom{\mathop{\mathcal{H}\!\mathit{om}}\nolimits}
\def\SheafExt{\mathop{\mathcal{E}\!\mathit{xt}}\nolimits}
\def\Sch{\mathit{Sch}}
\def\Mor{\mathop{\mathrm{Mor}}\nolimits}
\def\Ob{\mathop{\mathrm{Ob}}\nolimits}
\def\Sh{\mathop{\mathit{Sh}}\nolimits}
\def\NL{\mathop{N\!L}\nolimits}
\def\CH{\mathop{\mathrm{CH}}\nolimits}
\def\proetale{{pro\text{-}\acute{e}tale}}
\def\etale{{\acute{e}tale}}
\def\QCoh{\mathit{QCoh}}
\def\Ker{\mathop{\mathrm{Ker}}}
\def\Im{\mathop{\mathrm{Im}}}
\def\Coker{\mathop{\mathrm{Coker}}}
\def\Coim{\mathop{\mathrm{Coim}}}

% Boxtimes
%
\DeclareMathSymbol{\boxtimes}{\mathbin}{AMSa}{"02}

%
% Macros for moduli stacks/spaces
%
\def\QCohstack{\mathcal{QC}\!\mathit{oh}}
\def\Cohstack{\mathcal{C}\!\mathit{oh}}
\def\Spacesstack{\mathcal{S}\!\mathit{paces}}
\def\Quotfunctor{\mathrm{Quot}}
\def\Hilbfunctor{\mathrm{Hilb}}
\def\Curvesstack{\mathcal{C}\!\mathit{urves}}
\def\Polarizedstack{\mathcal{P}\!\mathit{olarized}}
\def\Complexesstack{\mathcal{C}\!\mathit{omplexes}}
% \Pic is the operator that assigns to X its picard group, usage \Pic(X)
% \Picardstack_{X/B} denotes the Picard stack of X over B
% \Picardfunctor_{X/B} denotes the Picard functor of X over B
\def\Pic{\mathop{\mathrm{Pic}}\nolimits}
\def\Picardstack{\mathcal{P}\!\mathit{ic}}
\def\Picardfunctor{\mathrm{Pic}}
\def\Deformationcategory{\mathcal{D}\!\mathit{ef}}
\usepackage{cleveref}
\crefname{originallemma}{original lemma}{original lemmas}
\Crefname{originallemma}{Original lemma}{Original lemmas}
\makeatletter
\newlabel{lemma-colimit}{{0A37}{source}{}{sourcepointer-1}{}}
\newlabel{lemma-colimit@cref}{{[originallemma][0][]0A37}{[1][0][]source}{}{}{}}
\makeatother
\crefname{originallemma}{original lemma}{original lemmas}
\Crefname{originallemma}{Original lemma}{Original lemmas}
\makeatletter
\newlabel{topology-lemma-make-spectral-space}{{0A31}{source}{}{sourcepointer-2}{}}
\newlabel{topology-lemma-make-spectral-space@cref}{{[originallemma][0][]0A31}{[1][0][]source}{}{}{}}
\makeatother
\crefname{originallemma}{original lemma}{original lemmas}
\Crefname{originallemma}{Original lemma}{Original lemmas}
\makeatletter
\newlabel{topology-lemma-make-spectral-space-minus}{{0AP0}{source}{}{sourcepointer-3}{}}
\newlabel{topology-lemma-make-spectral-space-minus@cref}{{[originallemma][0][]0AP0}{[1][0][]source}{}{}{}}
\makeatother
\crefname{originallemma}{original lemma}{original lemmas}
\Crefname{originallemma}{Original lemma}{Original lemmas}
\makeatletter
\newlabel{topology-lemma-spectral-sub}{{0902}{source}{}{sourcepointer-4}{}}
\newlabel{topology-lemma-spectral-sub@cref}{{[originallemma][0][]0902}{[1][0][]source}{}{}{}}
\makeatother
\title{Cohomology on the set of specialisations of a quasi-compact subset as a colimit over its neighbourhoods}
\author{The Stacks Project Authors}
\date{stacks-project:cohomology.tex:local-snapshot-20261011; GNU FDL 1.2 or later}
\begin{document}
\maketitle
\noindent\textit{Editorial note.} All delivered mathematical text is the author's own, copied from the pinned Stacks Project chapter file cohomology.tex; no statement or proof step was rewritten, reordered or removed.\par\medskip
\noindent\textit{Editorial note.} The second paragraph is the author's own 'The same proof works for part (2) except we use ...'; the sentence is reproduced verbatim and the variant result is carried as a navigation stub.\par\medskip
\noindent\textit{Editorial note.} The counted spans carry 4 navigation stubs and no bibliography citation.\par\medskip
\noindent\textit{Editorial note.} The preamble is the chapter's retained minimal preamble (sources/stacks-cohomology/retained-preamble.tex lines 1-81), a verbatim excerpt of the native preamble.\par\medskip
\noindent\textit{Editorial note.} The label lemma-cohomology-of-neighbourhoods-of-closed is defined only in cohomology.tex, so no cross-chapter restatement had to be weighed for this card.\par\medskip
\noindent\textit{Allowed standard prerequisites (editorial).}\begin{itemize}
\item Structure of the set of specialisations of a quasi-compact subset as an inverse limit of quasi-compact open neighbourhoods. Topology, Lemma topology-lemma-make-spectral-space, invoked at proof lines 3884-3885. Supplies the presentation of W as an inverse limit.
\item Subspaces appearing as quasi-compact opens are spectral. Topology, Lemma topology-lemma-spectral-sub, invoked at proof lines 3886-3887. Supplies that the colimit formula is applicable to each neighbourhood.
\item Colimit formula for cohomology on a spectral space presented as an inverse limit. Same chapter, Lemma lemma-colimit, invoked at proof lines 3888. Supplies both statements of the lemma.
\item The same structure theorem for the complement of a constructible subset. Topology, Lemma topology-lemma-make-spectral-space-minus, invoked at proof lines 3889-3890. Supplies the second statement.
\end{itemize}\medskip
\ifdefined\hypertarget\hypertarget{sourcepointer-1}{}\fi
\noindent\textit{Source reference.} Original lemma 0A37: sources/stacks-cohomology/pinned-original-source.tex, source lines 3360--3360. Pointer outside the counted span, invoked once for context; the load-bearing argument of the counted statement is the author's proof reproduced in full in the proof span.\par\medskip
\ifdefined\hypertarget\hypertarget{sourcepointer-2}{}\fi
\noindent\textit{Source reference.} Original lemma 0A31: sources/stacks-chow/tags.txt, source lines 12724--12724. Pointer outside the counted span, invoked once for context; the load-bearing argument of the counted statement is the author's proof reproduced in full in the proof span.\par\medskip
\ifdefined\hypertarget\hypertarget{sourcepointer-3}{}\fi
\noindent\textit{Source reference.} Original lemma 0AP0: sources/stacks-chow/tags.txt, source lines 13458--13458. Pointer outside the counted span, invoked once for context; the load-bearing argument of the counted statement is the author's proof reproduced in full in the proof span.\par\medskip
\ifdefined\hypertarget\hypertarget{sourcepointer-4}{}\fi
\noindent\textit{Source reference.} Original lemma 0902: sources/stacks-chow/tags.txt, source lines 11395--11395. Pointer outside the counted span, invoked once for context; the load-bearing argument of the counted statement is the author's proof reproduced in full in the proof span.\par\medskip
\begin{lemma}
\label{lemma-cohomology-of-neighbourhoods-of-closed}
Let $X$ be a spectral space. Let $\mathcal{F}$ be an abelian sheaf on $X$.
Let $E \subset X$ be a quasi-compact subset. Let $W \subset X$ be the set of
points of $X$ which specialize to a point of $E$.
\begin{enumerate}
\item $H^p(W, \mathcal{F}|_W) = \colim H^p(U, \mathcal{F})$
where the colimit is over quasi-compact open neighbourhoods of $E$,
\item $H^p(W \setminus E, \mathcal{F}|_{W \setminus E}) =
\colim H^p(U \setminus E, \mathcal{F}|_{U \setminus E})$
if $E$ is a constructible subset.
\end{enumerate}
\end{lemma}

\begin{proof}
From Topology, Lemma \ref{topology-lemma-make-spectral-space}
we see that $W = \lim U$ where the limit is over the quasi-compact
opens containing $E$. Each $U$ is a spectral space by
Topology, Lemma \ref{topology-lemma-spectral-sub}.
Thus we may apply Lemma \ref{lemma-colimit} to conclude that (1) holds.
The same proof works for part (2) except we use
Topology, Lemma \ref{topology-lemma-make-spectral-space-minus}.
\end{proof}

\end{document}
